\documentclass[11pt]{article}
\usepackage[margin=1in]{geometry}
\usepackage[T1]{fontenc}
\usepackage{xurl}
\setlength{\emergencystretch}{2em}
\title{Tactix: stable public Lucky actions across proof modes}
\author{Codex for Andrew}
\date{27 September 2026}
\begin{document}
\maketitle

\paragraph{Change and reason.}
The terminal Opus AM473 campaign found eight random public-program twins whose
first difference was Lucky enumerating 19 actions with explicit DSR output and
21 without proof. There was no SAT/UNSAT contradiction, but the proof writer
changed the sampled public action. Lucky now always enumerates stable IDs
0--20. Startup-PBP actions 19 and 20 still decline before mutation when the
original-input boundary or selected proof format cannot certify them. Proof
routing, public X7 legality and DefaultPolicy's explicit startup schedule are
unchanged.

\paragraph{Focused validation.}
A public C-API regression compares proofless and strict-DSR enumeration,
applies both unsupported startup actions and confirms an empty DSR stream.
Release and proof-shadow Debug API, parity/count and random-policy tests pass.
Both builds' parity and count startup PBPs independently verify against their
original CNFs through VeriPB and CakePB; the missing-derivation count control
is rejected. This is affected-scope evidence, not full N1/N2 acceptance.

\paragraph{Remaining work.}
Opus's 30 direct Proof API D2F flags (29 DSR, one DRAT) remain an open
production-reach question. Those scripts can bypass kernel root-state
maintenance; the public X7 validator checks live root units. Other producers
emit transient proof lemmas, so their checker-visible lifecycle needs a
targeted public reach audit before classifying D2F as out of contract. A new
Opus campaign and full Debug run on this source are due, followed by the
automatic-proof Rank/B28+C5, WV100 and WV2-89 gates. N1--N4 remain open; Tier
D has not been submitted.

\paragraph{Evidence.}
\url{benchmarks/reports/2026-09-27-am482-stable-lucky-public-enumeration.md};
\url{benchmarks/reports/2026-09-27-am481-opus-terminal-debug-checkpoint.md}.
\end{document}
